\section{真实研究迭代与证据核验}
\subsection{七轮研究的运行轨迹}
表~\ref{tab:research-turns} 根据原版事件、逐轮摘要及最终外部审计形成，覆盖 2026 年 10 月 25 日完成的七个研究轮次。单次研究轮次可包含多个子任务，也可能失败；不能将“四阶段菜单”理解为四次调用就完成研究。

\begin{table}[htbp]
\centering
\caption{七轮研究证据；时间为当日日本时间，调用列为接受数/失败数}
\label{tab:research-turns}
\small
\setlength{\tabcolsep}{4pt}
\begin{tabular}{clrrrrp{0.22\linewidth}}
\toprule
轮次 & 时间与主要任务 & 调用 & Python 成功 & 审核通过 & 归档 & 状态 \\
\midrule
1 & 15:44 建模校验 & 9/1 & 3 & 0 & 8 & 局部未完成 \\
2 & 16:02 补证复核 & 15/0 & 1 & 1 & 5 & 局部闭环完整 \\
3 & 16:23 POD 与失效分析 & 13/0 & 2 & 1 & 3 & 局部闭环完整 \\
4 & 16:47 改进首次尝试 & 2/1 & 0 & 0 & 4 & 模型请求超时 \\
5 & 16:56 改进内核开发 & 13/0 & 1 & 1 & 4 & 局部闭环完整 \\
6 & 17:31 冻结最终评价 & 13/0 & 7 & 0 & 7 & 实验完成，审核材料未充分 \\
7 & 17:53 热源补证复审 & 12/0 & 1 & 1 & 7 & 局部闭环完整 \\
\midrule
合计 & 全部已结束轮次 & 77/2 & 15 & 4 & 38 & 4 个局部闭环完整 \\
\bottomrule
\end{tabular}
\end{table}

“调用”是本地 Provider 的预约及响应记录，不是服务端重试次数或精确计费请求数。Python 成功表示原版工具实际返回退出码 0，不表示相同数量的新数学定理；共 16 次 Python 工具尝试中有 15 次成功。审核列统计整体审核通过的次数；归档条目可包含计划、失败和局部结论，不能仅凭存在归档判断研究完成。共 6 次整体审核尝试，其中 2 次因材料不足未定、4 次通过。38 条归档均与原版研究记忆匹配。

七轮记录的配置模型均为 \texttt{gpt-6-astra}。假设、计算、逻辑三个维度是该配置下具有不同审核职责的模型调用，不是三个互相独立的数学证明。文中的“外部独立复核”则指另外实现的组装、程序、推导和复算证据，两者不能混称。

七轮墙钟耗时之和为 $7407.187$ 秒，约 $123.45$ 分钟，包括模型等待、工具运行与原版审核归档，但不包括轮间协调、额外独立复核和报告编写的全部时间。它属于研究运行成本，与毫秒或秒量级的数值服务计时分开；不能据数值求解时间声称整项智能体研究仅耗时几秒。

\subsection{失败怎样改变后续行动}
\textbf{第一轮：运行产物存在，但接口与修订尝试失败。}
智能体实际写出求解器并运行校验，但尝试直接写入原版管理的 \texttt{workspace/problem.md} 被保护机制拒绝，没有绕过该限制。它随后误判序列化代码需要替换，修订脚本因找不到预期匹配而退出 1；检查失败本身保留。该轮最终还出现工具参数 JSON 中未转义控制字符，Provider 拒绝响应，因而没有通过整体审核。即使已有 3 次 Python 成功执行和 8 条归档，运行器仍将本轮标为不完整。

\textbf{第二轮：补充审核可见证据，保留未定结论。}
重新执行基础校验后，首次整体审核的计算维度因给定材料不足而未定。后续补充了装配实现、检查定义、容差与真实输出，第二次整体审核在 16:17:32 通过假设、计算和逻辑三个维度。后一次通过没有删除前一次未定；它只支持限定的离散命题与提交材料一致性，不是审核模型独立重算全部数组或完成形式化验证。

\textbf{第三轮：从调参达标推进到开发失效。}
原版工具分别执行固定 POD 基线与失效分析。智能体保存了调参选秩结果，又用 16 个角点发现两例场误差超标，并报告屏障指标的保守触发。该轮不是只生成一份研究计划，而是具有两个实际程序执行、一次整体审核和原版归档的局部闭环。屏障的外部先验来源在产物中明确标注。

\textbf{第四、五轮：将超时的大任务拆成可检查的内核。}
第四轮一次模型请求超过 300 秒后失败，没有 Python 实验执行；失败背景仍被归档。第五轮先集中实现参数分离、快速完整求解、固定降阶与自适应更新内核，再实际运行开发校验，653 项检查通过，随后完成局部审核与归档。该内核开发并不等于最终三网格、256 参数和完整成本评价已经完成；最终评价另按冻结协议执行与记账。

\subsection{如何核验，以及不能由此推出什么}
外部 \texttt{audit\_evidence.py} 只读取运行记录，不调用模型也不新增实验。它核对请求预算、工具调用与结果配对、成功退出码、审核返回值、原版归档条目、代码快照的 SHA-256，以及上游固定版本。脚本还检查 Codex Provider 的原生事件，确认接受的响应没有另外执行 Codex 自带工具；研究工具的执行仍经 Moonshine。这个来源审计证明的是可追踪性与记录一致性，不证明数学命题绝对正确。

数学与数值复核由另一组独立文件承担：独立装配、制造解、POD 子空间复算、角点误差复算和有理数屏障检查，与原生运行记录分别保留。局部闭环完整的判定同时要求实际程序执行、审核通过和归档，且没有运行失败；没有这些证据时，即使模型文字乐观，也仅记为部分进展。

\textbf{第六、七轮：保留完整实验，补齐审核的模型说明。}第六轮实际完成冻结、九次实验及汇总；计算与逻辑审核通过，但假设审核因未展开两个矩形热源的交叠积分而未定。第七轮直接检查式~\eqref{eq:source-overlap} 与原函数、快速载荷的一致性，三网格逐项差为零、总热量满足 $0.12$ 的 $10^{-14}$ 容差；没有新增 FOM/ROM 求解或测试参数。补齐证据后综合审核通过，并实际归档 7 条记录。归档前生成的文本曾写“无回执”，最终事件与原版记忆已确认之后的归档，审计按实际事件判定。

每轮目录及研究阶段快照保留失败与修订前状态，本地适配器的版本变化也按轮保存，不能用当前文件倒写旧执行历史。研究结束后另完成交付可移植性修复：无 Git 元数据时核验固定源码清单、直接载入原版包、相对定位证据、允许未配置模型时退出菜单或查看记录。相关安装与离线检查在独立环境完成，模型调用为零。这些是外部交付修复，不声称七轮历史研究使用了后来的启动器版本；原版 163 个源文件及冻结科学程序保持不变。

\begin{samepage}
\subsection{复算与继续研究}
交付目录的 \texttt{README.md} 给出安装和完整使用说明。Windows 下使用 Python 3.12 的 64 位版本运行“安装环境.cmd”，安装固定科学依赖；该步骤与离线检查不发起模型研究。以下命令在项目根目录运行：
\begin{verbatim}
.\.venv\Scripts\python.exe -B -X utf8 reproduce_thermal.py
.\.venv\Scripts\python.exe -B -X utf8 reproduce_thermal.py --full
\end{verbatim}
\end{samepage}
第一条在新副本中复算 $n=24$ 的一次重复，第二条复算全部九次；程序沿用冻结参数、核对科学结果与自适应轨迹，并检查原工作区未被改变。新耗时不要求与原观测相同。首次简短复算已经通过。若要继续生成研究方案，使用“启动热仿真研究.cmd”；这条路径需要用户自己的 Codex 登录与可用模型，会产生新的研究调用与记录。报告、既有数值复算和启动新的智能体研究是三个不同用途。
